% Inclusion fragment. Requires amsmath, amssymb, amsthm.
% Uses theorem, lemma, corollary, proof environments supplied by the parent.
\section{A general obstruction to computing the initializer from the update}
\label{sec:tied-seed-circuits}

The current counter-congruence and transient-seed drafts already prove a
one-coordinate majorization theorem for the particular operator
\[
 B(0;u,v.h)=1+h(0,0),\qquad
 B(n+1;u,v.h)=h(n,B(n;u,v.h)).
\]
The transient-seed draft also states composition closure and the
nonexistence of an injective pairing term. We do not claim those
conclusions again as new results. The following extension
identifies a substantially broader family of initializer-elimination
attempts to which the same obstruction applies. It allows finitely many
adaptive probes of the update and arbitrary primitive-recursive
combinations of the answers. The restriction is that the initializer
receives no independent unbounded numerical input.

\paragraph{Fixed query circuits.}
Fix a finite acyclic circuit $P$ with no ordinary input nodes. Its nodes
are numerals; applications of fixed primitive-recursive functions to
earlier node values; or queries $H(a,b)$, whose two arguments are earlier
node values. One node is designated as output. The circuit may query $H$
several times, and a later query may depend on earlier answers. Its finite
graph and all its ordinary primitive-recursive operations are fixed by
the constructor, rather than supplied as numerical inputs at runtime.

For a fixed update term $h(u,v,\vec p)$, let
$h_{\vec p}(a,b)=h(a,b,\vec p)$. Define the tied-seed recursor
\[
 \begin{split}
 R_P(0;u,v.h)&=P^{h_{\vec p}},\\
 R_P(n+1;u,v.h)&=h(n,R_P(n;u,v.h),\vec p).
 \end{split}
 \tag{Tied seed}
\]
The only access of the circuit to $\vec p$ is through queries to
$h_{\vec p}$. In particular it does not receive a separate initializer,
the actual counter argument, or a passive numerical parameter as an
ordinary input. Each fixed instance denotes a primitive-recursive
function: the circuit is finite and the outer recursion is ordinary
primitive recursion. No universal evaluator is introduced.

\paragraph{Permitted term primitives.}
The term language may contain constants, projections, and arbitrary fixed
unary primitive-recursive functions. More generally, it may contain any
fixed primitive-recursive operation $f$ for which either a constant bound
is known or there are a fixed input coordinate $j$ and unary
primitive-recursive function $F$ satisfying
\[
                    f(\vec x)\le F(x_j).
 \tag{One-coordinate bound}
\]
Thus predecessor, zero tests, bounded Boolean tests, minimum, and
one-sided data gates are all permitted. Ordinary finite substitution
and arbitrary nesting of the displayed $R_P$ constructors are allowed.
The arbitrary combining operations inside a query circuit need not obey
this bound; they receive only the circuit's internally generated values.

\begin{lemma}[A circuit cannot acquire an independent unbounded input]
For any fixed query circuit $P$ the following statements hold.
\begin{enumerate}
\item If $H(a,b)\le M$ for all $a,b$, then
      $P^H\le C_P(M)$ for some unary primitive-recursive $C_P$,
      independent of $H$.
\item If $H(a,b)\le F(a)$ for all $a,b$, where $F$ is a fixed unary
      primitive-recursive function, then $P^H$ has a constant bound
      depending on $P,F$, and independent of $H$. The same statement
      holds with $F(b)$ in place of $F(a)$.
\end{enumerate}
\end{lemma}
\begin{proof}
For a fixed $r$-ary primitive-recursive operation $g$, put
\[
 \widehat g(z)=\max\{g(z_1,\ldots,z_r):z_1,\ldots,z_r\le z\}.
\]
This is a nondecreasing unary primitive-recursive function. Nullary
operations are simply their constant values.

Traverse the finite acyclic circuit in a topological order. For the
first assertion, give each numeral its constant bound, each oracle
answer the bound $M$, and each ordinary operation node the bound
$\widehat g$ of the maximum of its predecessor bounds. Every node bound
is unary primitive recursive in $M$; the output bound is $C_P(M)$.

For the second assertion use constant bounds throughout the traversal.
At an oracle query, apply the nondecreasing envelope
$\widehat F(z)=\max_{t\le z}F(t)$ to the bound on the controlling query
argument. Since its argument node precedes the query, its constant bound
is already available. This gives a finite constant bound at every node.
The proof for the second query coordinate is identical.
\end{proof}

\begin{theorem}[Tied-seed circuit obstruction]
Every fixed term $t(x_1,\ldots,x_k)$ of the language just described has
either a constant bound or a bound
\[
              t(x_1,\ldots,x_k)\le F(x_i)
\]
for some one coordinate $i$ and unary primitive-recursive function $F$.
The coordinate and the function may depend on the term. This is an
external standard-function statement; it makes no claim about a uniform
primitive-recursive map from arbitrary terms to their bounding functions.
\end{theorem}
\begin{proof}
Replace every unary bounding function by its nondecreasing envelope.
The property is preserved by a primitive operation with a known bound:
only the term placed in its controlling input matters, and the two
nondecreasing unary bounds can be composed. Constants and projections
are immediate. It remains to check the new recursion constructor.

Apply the structural induction hypothesis to
$h(u,v,\vec p)$. Write $r_j(\vec p)$ for its recursive state at counter
$j$.

If $h(u,v,\vec p)\le F(p_i)$ for a passive parameter, the first circuit
lemma gives
\[
 r_0(\vec p)\le C_P(F(p_i)),\qquad
 r_{j+1}(\vec p)\le F(p_i).
\]
Consequently $\max(C_P(F(p_i)),F(p_i))$ bounds the result independently
of the actual counter. A constant bound for $h$ yields a constant bound
in exactly the same way.

If $h(u,v,\vec p)\le F(u)$, the second circuit lemma bounds the initial
state by a constant $b$. Every subsequent state satisfies
$r_{j+1}\le F(j)$. Therefore the nondecreasing unary primitive-recursive
function
\[
                       M(n)=\max(b,F(n))
\]
bounds $r_n$ at every $n$, including $n=0$.

If $h(u,v,\vec p)\le F(v)$, the second circuit lemma again gives a
constant bound $b$ on the initial state. Put
\[
                       K(z)=\max(z,F(z)).
\]
Then $K$ is nondecreasing and $K\ge\operatorname{id}$. Induction on $n$
gives $r_n\le K^n(b)$. Fixed-function iteration is primitive recursive,
and $n\mapsto K^n(b)$ is nondecreasing.

In the last two cases the actual counter is itself a smaller term.
Compose its constant or one-coordinate bound with $M$ or
$n\mapsto K^n(b)$, respectively. These cases exhaust the possible
controlling coordinates of $h$, completing the structural induction.
\end{proof}

\begin{corollary}[Combining two unbounded data inputs remains impossible]
This term language cannot define ordinary binary addition, maximum, a
general three-argument numerical selector, or an injection
$\mathbb N^2\longrightarrow\mathbb N$.
\end{corollary}
\begin{proof}
For addition or maximum, hold the asserted controlling input fixed and
increase the other input. For a selector $C(c,a,b)$ that returns $a$ at
$c=0$ and $b$ at $c>0$, each proposed controlling coordinate is defeated
by fixing it and increasing a selected data input. For an injection,
fixing the controlling coordinate leaves infinitely many argument pairs
whose outputs would all lie in one finite bounded interval. Constant
bounds fail for the same reasons.
\end{proof}

\paragraph{What this does and does not establish.}
The original particular $B$ operator is recovered by a one-query circuit
followed by successor. The extension also rules out taking several
fixed or adaptive probes of that same update, then adding, multiplying,
maximizing, or otherwise primitively-recursively combining their answers
to manufacture a seed. It continues to hold after adjoining arbitrary
fixed unary primitive-recursive functions and the other bounded-input
primitives described above. It is therefore stronger than checking a
single seed recipe or searching a finite collection of candidate terms.

This remains a standard-function \emph{term-definability} obstruction.
It does not rule out quantified graph definitions, a relational
interpretation, a different interpretation domain, additional axioms, or
bi-interpretability obtained by another route. An independent unbounded
initializer is one way to escape the hypothesis, but the theorem does
not assert that it is the only way. The theorem is not a smaller positive
axiomatization and is not presented as a globally novel result.

\paragraph{Why the interface restriction is necessary.}
If the initializer circuit is allowed one further independent passive
input $y$, take $h(u,v,x,y)=x$ and let the initializer be
$y+h(0,0,x,y)$. At counter zero the resulting term is exactly $x+y$,
so the one-coordinate conclusion fails. Likewise, a circuit permitted
two independently chosen update oracles can query the constant updates
$h_1=x$ and $h_2=y$ and add their answers. These observations identify
the boundary of the theorem; they do not claim that either relaxation
alone supplies a complete arithmetic basis or a smaller proof checker.
